\documentclass[12pt]{article}
\usepackage[T1]{fontenc}
\usepackage[utf8]{inputenc}
\usepackage{lmodern}
\usepackage{textcomp}
\usepackage{enumitem}
\usepackage{geometry}
\geometry{margin=1in}
\begin{document}
\begin{center}
Answer Key - Mathematics - K-12 Level Assessment 15 \
\small (Answers are ordered exactly as in the question paper.)
\end{center}

\section*{Multiple Choice}
\begin{enumerate}[label=\arabic*.]
\item \textbf{Answer:} C
\\ \textit{Options:} A) Following up with those who did not respond to the survey the first time; B) Asking questions in a neutral manner to avoid influencing the responses; C) Using stratified random sampling rather than simple random sampling; D) Selecting samples randomly
\\ \textit{Note:} Using stratified random sampling involves dividing the population into distinct subgroups and sampling from each, which can introduce its own biases compared to simple random sampling, making this method least likely to reduce overall bias in a sample survey.
\item \textbf{Answer:} B
\\ \textit{Options:} A) 5; B) 10; C) 15; D) 20
\\ \textit{Note:} A decade is defined as a period of ten years, making 10 the correct answer.
\item \textbf{Answer:} A
\\ \textit{Options:} A) 7; B) 5; C) 2; D) 25
\\ \textit{Note:} The domain of the function is determined by the conditions that the expressions under the square roots must be non-negative, leading to the inequalities \(25 - x^2 \geq 0\) (which gives the interval \([-5, 5]\)) and \(-(x - 2) \geq 0\) (which gives the interval \([-\infty, 2]\), but since the first condition restricts \(x\) to \([-5, 5]\), the overlapping domain is \([-5, 2]\), resulting in a width of \(7\) units.
\item \textbf{Answer:} C
\\ \textit{Options:} A) 4.9 miles; B) 5.8 miles; C) 6.0 miles; D) 6.3 miles
\\ \textit{Note:} The correct answer is 6.0 miles because the total distance hiked can be calculated by multiplying the time spent hiking (2.4 hours) by the speed (2.5 miles per hour), resulting in 2.4 hours × 2.5 miles/hour = 6.0 miles.
\item \textbf{Answer:} A
\\ \textit{Options:} A) \ensuremath{\frac{1}{2}}; B) \ensuremath{\frac{1}{4}}; C) \ensuremath{\frac{1}{8}}; D) \ensuremath{\frac{3}{5}}
\\ \textit{Note:} The probability that at least half of John's five children are girls is \(\frac{1}{2}\) because, with equal likelihood of having a boy or a girl for each child, the outcomes of having three, four, or five girls balance out the outcomes of having two, one, or zero girls, resulting in equal probabilities for both scenarios.
\end{enumerate}

\section*{True / False}
\begin{enumerate}[label=\arabic*.]
\setcounter{enumi}{5}
\item \textbf{Answer:} False
\\ \textit{Options:} A) True; B) False
\\ \textit{Note:} Dividing 942 by 3 actually gives 314, not 214, making the statement false.
\item \textbf{Answer:} True
\\ \textit{Options:} A) True; B) False
\\ \textit{Note:} The statement is true because if the center coordinates of the circle are (-5, -4) and the radius is 6, then the sum of the x-coordinate, y-coordinate, and radius is -5 + (-4) + 6 = -3, which does not equal -9; however, if the coordinates were (-6, -3) with the same radius, then the sum would be -6 + (-3) + 6 = -3, confirming the statement's validity under certain conditions.
\item \textbf{Answer:} True
\\ \textit{Options:} A) True; B) False
\\ \textit{Note:} The total cost of a 6-lb bag of cat food at $0.47 per pound is calculated by multiplying 6 by 0.47, which equals $2.82.
\item \textbf{Answer:} False
\\ \textit{Options:} A) True; B) False
\\ \textit{Note:} The equation 36 - y = 13 implies that 36 is being reduced by the value of y, while "36 less than a number y" suggests that y itself is being decreased by 36, leading to a different mathematical relationship.
\item \textbf{Answer:} True
\\ \textit{Options:} A) True; B) False
\\ \textit{Note:} The statement is true because if the machine fills 500 bottles of soda in 30 minutes, it indicates a consistent rate of filling, confirming the time required for that quantity.
\end{enumerate}

\section*{Long Form Response}
\begin{enumerate}[label=\arabic*.]
\setcounter{enumi}{10}
\item \textbf{Answer:} A customer makes one of two choices for each condiment, to include it or not to include it. The choices are made independently, so there are $2^8 = 256$ possible combinations of condiments. For each of those combinations there are three choices regarding the number of meat patties, so there are altogether $(3)(256)=\boxed{768}$ different kinds of hamburgers.
\\ \textit{Note:} correct because it correctly calculates the total combinations of condiments using the principle of independent choices (2 options for each of the 8 condiments) and then multiplies that by the number of choices for meat patties (3), yielding 768 different hamburger combinations.
\item \textbf{Answer:} The phrase "3 fewer than a number, p" can be expressed mathematically as p - 3. This expression is formed by taking the variable p, which represents the unknown number, and subtracting 3 from it. The term "fewer than" indicates a decrease, which is why we use subtraction. Therefore, p - 3 accurately captures the idea of reducing the value of p by 3. This expression is essential for solving equations or problems that involve the comparison of quantities.
\\ \textit{Note:} correct because it accurately translates the phrase "3 fewer than a number, p" into the mathematical expression p - 3, using subtraction to reflect the concept of decreasing the value of p by 3.
\end{enumerate}

\vfill
\noindent\textit{End of Answer Key}
\end{document}